% Verbatim copy of /Users/alingalatan/workspaces/grothe/handoff/manuscript/legacy/ch30-gpt-forest.tex (file dated 2026-10-25). Only this line was added.
% Consolidated source module (formerly Chapter 30). Owner: lower-bound/06-original-responses.tex
\section{Lower-bound GPT's reassessment: a diffuse theorem, the conditioning budget, and an adaptive-response saddle (20 October)}\label{ch:gpt-forest}
\chapsummary{Diffuse comparison, conditioning-budget obstructions and adaptive-response duals. Preserve the distinction between an arbitrary cubic law and a realizable original response.}

\paragraph{Provenance and status (IDE note, 20 October).} Lower-bound GPT's \code{proof/forest-addendum.tex}, ported verbatim apart from this header, local macros and the bibliography
entries, which were moved to the common bibliography. Package \pathref{gpt-reviews/forest-reassessment-20261020/}. \emph{Our replay} \st{C}: \code{verify\_diffuse\_constants.py} PASS
(exact Fraction arithmetic, 15 rational tensor cases) and \code{check\_inner\_bracket.py} PASS (256 residual identities, 512 minorant completions, 725 disk choices), both under
\code{python3 -I}. We re-read the proof of the diffuse theorem: Poincaré in the visible coordinates gives $r\le2I$, the payoff is $2$-Lipschitz, and the real two-dimensional
Malliavin--Stein bound with $\mathbb E\|\tau-C\|_{HS}^2\le\kappa/3$ and $\kappa\le40vI$ gives $d_W\le\sqrt{80I/3}$. The steps are sound \st{P}.
\providecommand{\EE}{\mathbb E}\providecommand{\CC}{\mathbb C}\providecommand{\RR}{\mathbb R}\providecommand{\HH}{\mathcal H}
\providecommand{\Rea}{\operatorname{Re}}\providecommand{\Ima}{\operatorname{Im}}\providecommand{\Var}{\operatorname{Var}}
\providecommand{\tr}{\operatorname{tr}}\providecommand{\phase}{\operatorname{phase}}\providecommand{\eps}{\varepsilon}

\subsection{What the remaining problem is}

The current manuscript records the certified interval
\[
1.401674103\stale{22 Oct}{now $\KG>1.401695322458$}<K_G^{\CC}<1.404599424\stale{22 Oct}{now $\KG<1.404598161566$}
\]
on master commit \texttt{20eb917}. Nothing in this addendum changes it.
The global coherent Wick A theorem settles a genuine shared-driver family,
but no theorem identifies this family with the maximizers among general
cubic tensors. The missing step is a joint comparison or localization
theorem, not simply more accuracy on the already controlled family.

Two distinctions matter. First, a small coordinate coefficient in a
chosen basis is not the same as small influence in every direction.
Second, Conjecture K bounds every cubic in a relaxation, whereas the
original norm permits additional response fields to readjust to that
cubic. The latter distinction was already recognized by the retained
projection squares in manuscript \autoref{ch:weakk}. Section~\ref{forest:saddle}
gives a simpler exact organization of those squares.

\subsection{An explicit diffuse theorem}
\label{forest:diffuse}

Let \(x_1,\ldots,x_{2n}\) be independent standard real Gaussians, with
\(Z_j=(x_{2j-1}+ix_{2j})/\sqrt2\). For
\(Q\in\HH_{2,1}(\CC^n)\), define the real positive semidefinite matrix
\[
M_Q(i,j)=\EE\Rea\big[(\partial_iQ)\overline{\partial_jQ}\big],
\qquad
I(Q)=\lambda_{\max}(M_Q)
 =\sup_{\|h\|_{\RR^{2n}}=1}\EE|D_hQ|^2.
\]
All derivatives and eigenvalues in this definition use real coordinates.
In particular \(\tr M_Q=3\EE|Q|^2\).

For any integrable complex \(L=L(Z_1)\) and \(\eps>0\), set
\[
\Lambda(Q)=\EE(|L+Q|+|L-Q|)-\frac{\EE|Q|^2}{\eps},
\qquad
S=\sup_{v\ge0}\left(2\EE|L+\sqrt vG|-\frac v\eps\right),
\]
where \(G\) is an independent standard proper complex Gaussian.
The quantity \(S\) is finite, since its objective is at most
\(2\EE|L|+2\sqrt v-v/\eps\).

\begin{theorem}[Explicit approximate K in the diffuse regime]
\label{forest:thm-diffuse}
In every finite dimension, for every \(Q\in\HH_{2,1}\),
\[
\Lambda(Q)\le S+
\left(2\sqrt2+2\sqrt{\frac{80}{3}}\right)\sqrt{I(Q)}
\le S+14\sqrt{I(Q)}.
\]
No independence between \(Q\) and the visible \(Z_1\) is assumed.
\end{theorem}

\begin{lemma}[Real cubic estimates]
\label{forest:realcubic}
If \(F=I_3(f)\) is a real cubic chaos of variance \(v_F\), then
\[
\kappa_4(F)\le30v_F I(F),\qquad
\Var\|DF\|^2\le2\kappa_4(F).
\]
Here \(I(F)=\lambda_{\max}(\EE[DF(DF)^T])\).
\end{lemma}
\begin{proof}
Use the multiple-integral normalization \(v_F=6\|f\|^2\), and put
\[
A=\|f\otimes_2 f\|^2=\|f\otimes_1 f\|^2,\qquad
B=\|f\widetilde\otimes_1 f\|^2\le A.
\]
The equality of the two unsymmetrized contraction norms follows by
reordering their four tensor sums. The cubic product formula gives
the squared norms of the chaos \(0,2,4,6\) parts of \(F^2\) as
\[
36\|f\|^4,\quad 648A,\quad1944B,\quad
72\|f\|^4+648A.
\]
For the last expression one may use explicitly
\[
\|\operatorname{sym}(f\otimes f)\|^2
 =\binom63^{-1}\sum_{r=0}^3\binom3r^2
       \|f\otimes_r f\|^2
 =\frac{2\|f\|^4+18A}{20}.
\]
Thus
\[
\kappa_4(F)=1296A+1944B.
\]
The identity \(\|DF\|^2=(\mathcal L+6)F^2/2\), where
\(\mathcal L\) has eigenvalue \(-k\) on chaos \(k\), also gives
\[
\Var\|DF\|^2=2592A+1944B\le2\kappa_4(F).
\]
Finally \(\EE[DF(DF)^T]=18(f\otimes_2 f)\). The latter contraction
is positive semidefinite with trace \(\|f\|^2\), so
\[
A\le \frac{I(F)}{18}\frac{v_F}{6}.
\]
Since \(1296+1944=3240\), the first claim follows. Its constant is
sharp for \(F=\sqrt{v_F}H_3(x)/\sqrt6\).
\end{proof}

\begin{lemma}[Circular cubic Gaussian approximation]
\label{forest:clt}
Let \(q\in\HH_{2,1}\), \(v=\EE|q|^2>0\), and
\(\kappa=\EE|q|^4-2v^2\). In the Euclidean Wasserstein distance on
\(\CC=\RR^2\),
\[
\kappa\le40v I(q),\qquad
d_W(q,\sqrt vG)\le\sqrt{\frac{2\kappa}{3v}}
                 \le\sqrt{\frac{80}{3}I(q)}.
\]
\end{lemma}
\begin{proof}
The global phase action implies circularity of \(q\). Hence
\[
\Var(\Rea q)=v/2,\qquad
\kappa_4(\Rea q)=3\kappa/8.
\]
Lemma~\ref{forest:realcubic}, and \(I(\Rea q)\le I(q)\), prove the
first bound.

Write \(F=(u,w)=(\Rea q,\Ima q)\), \(C=(v/2)I_2\), and
\(\tau_{ij}=\langle DF_i,DF_j\rangle/3\). Define
\[
A_0=\|Du\|^2,\quad B_0=\|Dw\|^2,\quad C_0=\langle Du,Dw\rangle,
\quad
T_0=\Var A_0+\Var B_0+2\Var C_0.
\]
For \(F_\theta=u\cos\theta+w\sin\theta\), circularity and
Lemma~\ref{forest:realcubic} give
\(\Var\|DF_\theta\|^2\le3\kappa/4\). Direct angular integration yields
\[
\EE_\theta\Var\|DF_\theta\|^2
 =T_0/4+\Var(A_0+B_0)/8\ge T_0/4.
\]
Consequently \(\EE\|\tau-C\|_{\rm HS}^2=T_0/9\le\kappa/3\).
The multivariate Malliavin--Stein inequality of
Nourdin--Peccati--R\'eveillac, Theorem 3.5 and Lemma 3.3, states
\[
d_W(F,N(0,C))
\le \|C^{-1}\|_{\rm op}\|C\|_{\rm op}^{1/2}
       \big(\EE\|\tau-C\|_{\rm HS}^2\big)^{1/2}.
\]
Its matrix factor here is \(\sqrt{2/v}\), giving the result.
\end{proof}

\begin{proof}[Proof of Theorem~\ref{forest:thm-diffuse}]
Let \(\mathcal W=\sigma(Z_2,\ldots,Z_n)\),
\(q=\EE[Q\mid\mathcal W]\), \(R=Q-q\), and \(r=\EE|R|^2\).
The polynomial \(q\) remains in \(\HH_{2,1}\) and is independent of
\(Z_1\). Conditional Gaussian Poincar\'e gives
\[
r\le\EE(|\partial_1Q|^2+|\partial_2Q|^2)\le2I(Q).
\]
For every remote direction \(h\),
\(D_hq=\EE[D_hQ\mid\mathcal W]\). Conditional Jensen, together with
the vanishing visible derivatives of \(q\), gives \(I(q)\le I(Q)\).
By orthogonality and the 2-Lipschitz property of the norm sum,
\[
\Lambda(Q)\le\Lambda(q)+2\sqrt r-r/\eps
             \le\Lambda(q)+2\sqrt{2I(Q)}.
\]
The function
\(u\mapsto\EE_{Z_1}(|L+u|+|L-u|)\) is 2-Lipschitz.
Lemma~\ref{forest:clt} therefore gives
\(\Lambda(q)\le S+2\sqrt{80I(q)/3}\).
If \(\EE|q|^2=0\), the same conclusion holds directly. Combining
these bounds proves the theorem.
\end{proof}

\begin{corollary}[Explicit localization of a positive violation]
If \(\Lambda(Q)>S+\eta\), then \(I(Q)>(\eta/14)^2\).
At \(\eta=10^{-4}\), the cutoff is
\[
\frac{1}{19\,600\,000\,000}.
\]
\end{corollary}

This cutoff is structurally useful but computationally far too small.
Since \(\tr M_Q=3v\), the number of eigenvalues above a cutoff
\(\delta\) is at most \(3v/\delta\). Even at \(v<1\), the displayed
cutoff allows almost \(5.88\cdot10^{10}\) real heavy directions.
It does not produce a feasible finite cover.

For calibration, if
\[
Q_N=\sqrt{\frac vN}\sum_{j=1}^N H(X_j),\qquad
H(X)=\frac{X(|X|^2-2)}{\sqrt2},
\]
then \(I(Q_N)=3v/(2N)\) and
\(\EE|Q_N|^4-2v^2=42v^2/N\).
In contrast \(H(N^{-1/2}\sum_jX_j)\) has coordinate influences
of size \(3/(2N)\) but maximal directional influence \(3/2\).
Diffuseness must be measured invariantly, not in one chosen basis.

\subsection{Conditioning: the missing budget, exactly}
\label{forest:conditioning}

Fix the peak profile
\[
\eps=\frac{1891}{2500},\quad
L=Z(a+b|Z|^2),\quad a=\frac{699}{1250},\quad b=\frac{179}{1250}.
\]
Let \(q\) be a remote cubic and rotate an A-sector driver into one
coordinate \(X\):
\[
Q=q(w)+\alpha(|Z|^2-1)X+gG.
\]
Condition on \(Z,X,G\). The Wick expansion gives
\[
Q_P=\EE[Q\mid Z,X,G]
   =cH(X)+\alpha(|Z|^2-1)X+gG.
\]
This is the previously certified coherent family (a rotation of \(X\)
absorbs the phase of \(c\)). With \(R=Q-Q_P\) and
\(\Psi_L(u)=|L+u|+|L-u|\), define
\[
J_P(Q)=\EE\Psi_L(Q)-\EE\Psi_L(Q_P)\ge0.
\]
Conditional Jensen and orthogonality give the exact identity
\[
\Lambda(Q)=\Lambda(Q_P)+J_P(Q)-\|R\|_2^2/\eps.
\]
Thus the missing lifting estimate is precisely
\begin{equation}\label{forest:lift}
J_P(Q)\le\|R\|_2^2/\eps+[S-\Lambda(Q_P)]+\eta.
\end{equation}
For this reduction, \eqref{forest:lift} is equivalent to A-only
\(K_\eta\). It is not a proved estimate for arbitrary \(q\).

\begin{proposition}[The retained deficit cannot simply be removed]
The universal inequality \(J_P(Q)\le\|R\|_2^2/\eps\) is false,
even when \(Q\) is a genuine remote cubic.
\end{proposition}
\begin{proof}
For \(f(v)=\Lambda(\sqrt vG)\), Gaussian heat differentiation gives
\[
f'(0+)=\tfrac12\EE|L|^{-1}-1/\eps.
\]
The inverse moment is finite, and gamma-distribution Jensen yields
\[
\EE|L|^{-1}
=\int_0^\infty\frac{t^{-1/2}e^{-t}}{a+bt}\,dt
\ge\frac{\sqrt\pi}{a+b/2}.
\]
Therefore \(f'(0+)\ge\sqrt\pi/(2a+b)-1/\eps>0\).
The strict sign follows already from \(\pi>3\), since
\(3\cdot1891^2-3154^2=779927>0\).
Choose sufficiently small \(v>0\), and independent remote
\(Y_j\) unrelated to \(X\). Then
\[
q_m=\sqrt{v/m}\sum_{j=1}^mH(Y_j)
 \ \Longrightarrow\ \sqrt vG .
\]
Bounded second moments give uniform integrability of the norm
payoff; hence \(\Lambda(q_m)\to f(v)>f(0)\).
But \(\EE[q_m\mid Z,X]=0\).
The claimed stronger inequality would imply
\(\Lambda(q_m)\le\Lambda(0)=f(0)\), a contradiction for large finite \(m\).
\end{proof}

The same CLT argument shows
\(\sup_{n,q\in\HH_{2,1}(\CC^n)}\Lambda(q)\ge S\):
the benchmark is a diffuse boundary law. No monotonicity in \(m\)
is claimed. A finite-rank extremizer classification alone could miss
this boundary.

The book's tensor representation also rules out a tempting shortcut.
The cubic \(X^2\bar Y/\sqrt2\) has variance one, fourth cumulant 10,
and complex total-influence eigenvalues \(2,1\).
An independent Wick sum \(\sum c_jH(X_j)\) has corresponding
eigenvalues \(3|c_j|^2\) and cumulant \(42\sum|c_j|^4\).
Matching the eigenvalues would force cumulant \(70/3\), not 10.
Thus unitary rotation does not reduce general cubics to independent
Wick drivers. Here complex total influence uses
\(\EE|\partial_hq|^2+\EE|\bar\partial_hq|^2\);
it is a different normalization from the real matrix in
Section~\ref{forest:diffuse}.

One exact interpolation retains the missing joint information.
Let \(Q_t\) be conditional OU smoothing in the discarded coordinates
only, and let \(C_t=\sum_{j\perp}D_jQ_t(D_jQ_t)^T\), interpreting
complex values as real two-vectors. Then
\[
\Lambda(Q)-\Lambda(Q_P)
=\int_0^\infty\EE\tr\left[
\left(D_u^2\Psi_L(Q_t)-\frac2\eps I_2\right)C_t\right]dt.
\]
For a literal nonsingular formulation, replace each norm by
\(\sqrt{|\,\cdot\,|^2+\delta^2}\), prove the identity, and interpret
the displayed integral as the limit as \(\delta\downarrow0\).
Its value must be paid by the retained deficit in
\eqref{forest:lift}; its integrand is not universally nonpositive.

\subsection{An exact adaptive-response saddle}
\label{forest:saddle}

In fixed finite Gaussian dimension, let \(P=P_{1,0}\), \(Q=P_{2,1}\),
\(T=P-\eps Q\), and
\[
N_n(\eps)=\sup_{|f|,|g|\le1}\Rea\langle f,Tg\rangle.
\]
For \(U,x\in\operatorname{ran}P\), \(V,y\in\operatorname{ran}Q\), put
\begin{align*}
\Xi(U,V;x,y)
={}&\EE\bigl(|U+V-x-y|+|U-V+x-y|\bigr)\\
 &-\|U\|_2^2-\|V\|_2^2/\eps
   +\|x\|_2^2+\|y\|_2^2/\eps .
\end{align*}

\begin{theorem}[Two-field saddle]
\label{forest:thm-saddle}
For each fixed finite dimension,
\[
N_n(\eps)=\max_{U,V}\min_{x,y}\Xi(U,V;x,y).
\]
The inner minimizer exists and is unique. One can restrict the outer
maximum, after simultaneous unitary rotation, to \(U=tZ\) with
\[
0\le t\le\sqrt\pi/2,\qquad
\|V\|_2^2\le\eps^2(1-t^2).
\]
\end{theorem}
\begin{proof}
Put \(h=(f+g)/2\), \(k=(f-g)/2\). Exactly,
\[
\Rea\langle f,Tg\rangle
=\|Ph\|_2^2+\eps\|Qk\|_2^2-\|Pk\|_2^2-\eps\|Qh\|_2^2.
\]
The admissible \((h,k)\) satisfy \(|h+k|\le1\) and \(|h-k|\le1\).
Linearize the positive squares by maximizing
\[
2\Rea\langle U,h\rangle+2\Rea\langle V,k\rangle
-\|U\|^2-\|V\|^2/\eps.
\]
Dualize the negative squares by minimizing
\[
\|x\|^2+\|y\|^2/\eps
-2\Rea\langle x,k\rangle-2\Rea\langle y,h\rangle.
\]
For fixed \(U,V\), the response moment body is compact and convex,
and the latter variables enter continuously and convexly. The
finite-dimensional convex minimax theorem swaps this inner
minimization with the response maximization. There is no swap of
the outer maximization and the inner minimization. Pointwise response
maximization gives the two norms in \(\Xi\).

The inner objective is coercive and strictly convex. Its value
is at most \(2\|U\|+2\|V\|-\|U\|^2-\|V\|^2/\eps\), which tends
to \(-\infty\) as the outer norm tends to infinity. Continuity,
equivalently the compact response-moment representation, gives an
outer maximum. A maximizing original response pair supplies
\(U=Ph,V=\eps Qk\).
Now \(|h|\le1\) gives \(\|Ph\|\le\EE|Z|=\sqrt\pi/2\);
\(|h|^2+|k|^2\le1\) gives the variance budget.
A unitary rotation puts \(U\) in the displayed form.
\end{proof}

The supremum over \(n\) is still required for the full problem.
The theorem does not assert attainment in arbitrary dimension or
a bound on the number of influential coordinates.

\begin{proposition}[A posteriori bracket without singular derivatives]
\label{forest:bracket}
At any trial \(U,V,x,y\), choose
\[
f=\phase(U+V-x-y),\quad
g=\phase(U-V+x-y),\quad h=(f+g)/2,\quad k=(f-g)/2.
\]
At a zero argument choose any disk-valued support maximizer. Define
\[
R_{\rm in}=\|x-Pk\|^2+\|y-\eps Qh\|^2/\eps,\qquad
R_{\rm out}=\|U-Ph\|^2+\|V-\eps Qk\|^2/\eps.
\]
Then, with \(U,V\) fixed,
\[
\Xi-R_{\rm in}\le\min_{x',y'}\Xi(U,V;x',y')\le\Xi,
\]
and exactly
\[
\Xi=\Rea\langle f,Tg\rangle+R_{\rm in}-R_{\rm out}.
\]
\end{proposition}
\begin{proof}
For the chosen support responses,
the two norm terms equal
\(2\Rea\langle U-y,h\rangle+2\Rea\langle V-x,k\rangle\).
Completing the inner squares subtracts exactly \(R_{\rm in}\)
and leaves the concave response functional at this \(h,k\).
Its maximum is the inner minimum from the preceding proof,
so it is a lower bound for that minimum. Completing the remaining
outer squares gives the second identity.
\end{proof}

This bracket certifies accuracy of a numerical inner optimization
without a reciprocal-modulus Hessian. It does not certify a global
outer maximum. A sufficient uniform norm certificate would choose
explicit fields \(x(U,V),y(U,V)\) and prove
\(\Xi(U,V;x(U,V),y(U,V))\le B\) throughout the outer domain.
The fields may adapt to the complete tensor; general K is not
required for such a certificate.

At a smooth maximizing saddle the equations are
\[
U=Ph,\quad V=\eps Qk,\quad x=Pk,\quad y=\eps Qh.
\]
Thus \(f=\phase(Tg)\) and \(g=\phase(Tf)\).
With \(u=(U,V)\), \(v=(x,y)\), the necessary outer stability condition is
\[
H_{uu}-H_{uv}H_{vv}^{-1}H_{vu}\preceq0.
\]
The Hessian of \(\Xi\), when the integrals exist, is
\[
\EE\sum_{\pm}\frac{\Ima(\overline{w_\pm}d_\pm)^2}{|w_\pm|^3}
-2\|\delta U\|^2-2\|\delta V\|^2/\eps
+2\|\delta x\|^2+2\|\delta y\|^2/\eps,
\]
where \(w_\pm=U\pm V\mp x-y\) and
\(d_\pm=\delta U\pm\delta V\mp\delta x-\delta y\).
One must not use this formula at singularities without justification.
An alternative, valid for bounded real phase variations
\(\theta,\phi\) at a unimodular original maximizer, is
\[
\EE(|Tg|\theta^2+|Tf|\phi^2)
-2\Rea\langle i\theta f,T(i\phi g)\rangle\ge0.
\]
Neither stability condition presently implies the needed structural
classification.

\subsection{What one side can teach the other}

For an actual real response coefficient mixture \(c=(c_0,c_1,\ldots)\),
write \(D(c)=c_0-\sum_{k\ge1}|c_k|\). Manuscript \autoref{ch:crossbound}'s gap
identity for the cubic game is
\[
N_\eps-D(c)
=\bigl[N_\eps-(c_0-\eps c_1)\bigr]
 +(|c_1|-\eps c_1)+\sum_{k\ge2}|c_k|.
\]
Each term is nonnegative. For \(0<\eps<1\), a gap at most \(\delta\)
therefore forces
\[
|c_1|\le\frac{\delta}{1-\eps},\quad
\sum_{k\ge2}|c_k|\le\delta,\quad
N_\eps-(c_0-\eps c_1)\le\delta.
\]
The lower norm proof, cubic cancellation, and full higher-chaos tail
are distinct obligations. A common response pool can diagnose them
together; it cannot discard the tail or infer its sign.
The manuscript's approximately \(1.40226\) cubic ceiling is a
floating-point model value, not a certified endpoint. Proving K
would be valuable, but would not by itself reach the Haagerup value.

\subsection{Recommended change of target and verification status}

The priority is the original norm with adaptive response fields,
and the integrated lifting budget \eqref{forest:lift} when a coherent
core is retained. The genuinely diffuse part is now controlled by an
explicit theorem, but the generic constants are unusable at the
desired precision. Marginal Gaussian convergence cannot replace the
joint tensor comparison with the retained coordinates.

The proposed first experiment is to optimize the inner fields in
Theorem~\ref{forest:thm-saddle}, using Proposition~\ref{forest:bracket}
to distinguish solver error from a real obstruction, and inspect the
response-adjusted stability matrix on coupled tensors. A positive
outer direction rules out a proposed local maximizer; a nonpositive
finite matrix is only evidence and does not certify all directions.
This is a more targeted test than another cover of raw \(\Lambda\).

The accompanying exact-arithmetic script audits the tensor constants,
the circular Wick example, angular factors and cutoffs. These are
supporting checks, not formal proofs of the cited probability
theorems or of the infinite-dimensional optimization.
The live manuscript and prior certificate inputs are unchanged.

\subsection{IDE reading note}\label{sec:gpt-forest-reading}
Points of agreement with \autoref{ch:forest}, reached independently: no universal or pointwise principle is available; sector-by-sector or independent peeling fails; what is open is a
\emph{joint} statement about the coherent core and its remainder, with the deficit of the core retained. The two chapters complement each other. GPT's diffuse theorem makes the
\emph{near-Gaussian} regime rigorous for all sectors at once, but with a constant ($14\sqrt I$) far too weak for a certificate. Our \autoref{ch:forest} identifies which constants
should replace pointwise ones in that regime: averaged profile curvatures, $d_A\approx0.466$ against $0.292$. It also shows that the Gaussian-profile layer is closed, so a sharper diffuse
theorem has to target the third- and fourth-order joint cumulants, consistent with the $O(\kappa)$ rate of Chen--Chen--Liu~\cite{ccl}.
